\section{Vấn đề}
\label{sec:van-de}

\subsection{Bối cảnh}

Các mô hình ngôn ngữ lớn (LLM) đã vượt qua vai trò công cụ hỏi--đáp: khi được
ghép với khả năng gọi công cụ (tool calling), chúng có thể \emph{hành động} ---
tìm kiếm thông tin, đọc trang web, ghi dữ liệu, lên lịch --- thay vì chỉ sinh
văn bản. Lớp hệ thống này thường được gọi là \emph{AI agent}: LLM đóng vai trò
bộ điều khiển, quyết định bước tiếp theo dựa trên yêu cầu của người dùng và
quan sát thu được từ môi trường~\cite{yao2023react}.

Tuy nhiên, khoảng cách giữa một bản demo agent và một hệ thống dùng được hằng
ngày là rất lớn. Ba vấn đề nổi bật:

\begin{enumerate}
  \item \textbf{Độ tin cậy.} LLM không tất định và có thể ``ảo giác'': chọn sai
        công cụ, bịa tên công cụ không tồn tại, lặp vô hạn, hoặc thực hiện hành
        động có hại (xóa dữ liệu) mà người dùng không hay biết. Một agent hữu
        dụng phải có các lớp chặn tường minh: giới hạn số bước, timeout, xác
        nhận của con người trước hành động ghi/xóa.
  \item \textbf{Khả năng quan sát.} Khi agent trả lời sai hoặc hành xử lạ, câu
        hỏi ``nó vừa làm gì và vì sao'' phải trả lời được. Đa số demo agent chỉ
        ghi log rời rạc; việc dựng lại toàn bộ chuỗi quyết định (prompt nào,
        công cụ nào, tham số gì, quan sát gì, tốn bao nhiêu token) thường là
        không thể.
  \item \textbf{Đánh giá định lượng.} Hành vi agent phụ thuộc mô hình, prompt
        và chiến lược điều phối. Không có bộ đo lặp lại được thì mọi thay đổi
        (đổi model, sửa prompt) đều là cảm tính --- không biết hệ thống tốt lên
        hay xấu đi.
\end{enumerate}

\subsection{Bài toán}

Đồ án đặt bài toán: xây dựng một \textbf{AI agent cá nhân} giao tiếp qua
Telegram, tự thực hiện trọn vẹn vòng lặp \emph{nhận yêu cầu $\to$ phân tích
$\to$ lập kế hoạch $\to$ gọi công cụ $\to$ quan sát $\to$ điều chỉnh $\to$ trả
lời} trên một bộ công cụ thực dụng (tìm kiếm web, đọc trang, ghi chú, việc cần
làm, viết báo cáo, lên lịch định kỳ), với ba yêu cầu xuyên suốt tương ứng ba
vấn đề trên:

\begin{itemize}
  \item \textbf{An toàn theo lớp}: giới hạn bước và timeout ở mọi tầng; hành
        động ghi/xóa phải qua xác nhận người dùng (human-in-the-loop) với cơ
        chế tạm dừng/tiếp tục bền vững qua restart.
  \item \textbf{Trace là công dân hạng nhất}: mọi lệnh gọi LLM và mọi lần chạy
        công cụ đều được ghi vào cơ sở dữ liệu; một trace viewer dựng lại được
        toàn bộ timeline kèm token, chi phí, độ trễ.
  \item \textbf{Đánh giá tự động}: một eval harness chạy bộ test case cố định
        qua chính agent loop production, chấm tự động và xuất bảng metric ---
        làm nền cho thí nghiệm so sánh chiến lược và mô hình.
\end{itemize}

Tên dự án lấy ẩn dụ từ thí nghiệm tư duy của Pierre-Simon Laplace về một thực
thể biết toàn bộ trạng thái hiện tại của vũ trụ và từ đó suy ra tương lai ---
ở đây tượng trưng cho nguyên tắc thiết kế: agent \emph{quyết định dựa trên
trạng thái quan sát được và được ghi lại đầy đủ}, chứ không phải tuyên bố hệ
thống ``biết mọi thứ''.

\subsection{Phạm vi và đóng góp}

Phạm vi MVP: một người dùng cá nhân (hoặc nhóm nhỏ), một tiến trình, kênh
Telegram, 6 công cụ, SQLite. Ngoài phạm vi: multi-agent, RAG/vector memory,
hạ tầng phân tán.

Đóng góp chính của đồ án:

\begin{enumerate}
  \item Một agent loop \textbf{tự xây dựng} (không dùng framework agent) dưới
        dạng state machine tường minh, hỗ trợ hai chiến lược điều phối cắm
        chung một interface: ReAct~\cite{yao2023react} và
        Plan-and-Execute~\cite{wang2023planandsolve};
  \item Cơ chế human-in-the-loop \textbf{bền vững}: trạng thái tạm dừng lưu
        trong DB, tiếp tục được ở tiến trình khác, chống thực thi kép bằng
        compare-and-set;
  \item Hệ thống trace đầy đủ + trace viewer, dùng chung một nguồn dữ liệu cho
        cả gỡ lỗi lẫn đánh giá;
  \item Eval harness với 37 test case thuộc 8 nhóm (bao gồm cả prompt
        injection và phục hồi lỗi), 9 metric, tùy chọn
        LLM-as-judge~\cite{zheng2023judge}; và thí nghiệm so sánh
        $2 \times 2$ (chiến lược $\times$ mô hình).
\end{enumerate}
